\section{Simultaneous global arithmetic normalization}
\label{sec:simultaneous-global-cm}

\begin{theorem}
\label{thm:simultaneous-global-cm-normalization}
Let \(\Pi_{\rm global}\) be the set of primes \(p\geq5\) such that
\(Q_p(1)\) is odd, the natural first-layer map
\eqref{eq:uniform-natural-layer} is a finite isomorphism, and
\(\mathscr K_p\) is injective. Define
\[
 R_{\rm global}=R_{\rm quad}\cap
       \bigcap_{p\in\Pi_{\rm global}}\ker h_p,\qquad
 V_{\rm global}=V_{\rm quad}\cap R_{\rm global}.
\]
Then the joint images of \(R_{\rm quad}\), \(V_{\rm quad}\),
\(R_{\rm quad}\cap H\), and \(V_{\rm quad}\cap H\) are all
\[
 \mathcal P=\prod_{p\in\Pi_{\rm global}}P_p,
\]
including every finite sign subgroup. Consequently,
\begin{equation}
 R_{\rm quad}/R_{\rm global}\simeq\mathcal P,\qquad
 V_{\rm quad}/V_{\rm global}\simeq\mathcal P,\qquad
 G=R_{\rm global}H,
 \label{eq:uniform-global-factorization}
\end{equation}
and inclusion induces an isomorphism
\begin{equation}
 R_{\rm global}/V_{\rm global}
       \xrightarrow{\ \sim\ }R_{\rm quad}/V_{\rm quad}
 \label{eq:uniform-proodd-identification}
\end{equation}
which preserves the image of the arithmetic intersections.
In particular the quotient in \eqref{eq:uniform-proodd-identification}
is pro-odd. The relative quotients satisfy
\begin{equation}
 B(R_{\rm global})\simeq B,\qquad
 B(V_{\rm global})_{\max2}\simeq B_{\max2}.
 \label{eq:uniform-relative-quotients}
\end{equation}
The original restrictions of \(X_2\) to both new groups have
kernel exactly \(\mathcal S\), and image all continuous
\(G\)-invariant quadratic characters which vanish on the
corresponding arithmetic intersection.

The older unit set \(\Pi\) and the primes \(29,113\) belong to
\(\Pi_{\rm global}\), whereas \(163\) does not. No infinitude
assertion about \(\Pi_{\rm global}\) is needed.
\end{theorem}
\begin{proof}
For each \(p\in\Pi_{\rm global}\),
Theorem~\ref{thm:uniform-arithmetic-cm-image} gives the complete
\(P_p\)-axis inside the joint image of \(V_{\rm quad}\cap H\):
the witness is \(1\) in every other prime coordinate.
Finite products of such axes are dense in \(\mathcal P\).
The image of the compact group \(V_{\rm quad}\cap H\) is
closed, hence equals \(\mathcal P\). All four stated images
therefore equal \(\mathcal P\).

Matching every scalar coordinate by an element of
\(V_{\rm quad}\cap H\) gives
\[
 R_{\rm quad}=R_{\rm global}(V_{\rm quad}\cap H).
\]
The known factorization \(G=R_{\rm quad}H\) now gives
\eqref{eq:uniform-global-factorization}. The same equality,
and \(V_{\rm quad}\cap R_{\rm global}=V_{\rm global}\),
prove \eqref{eq:uniform-proodd-identification}.
If \(r\in R_{\rm quad}\cap H\), match its coordinates by
\(v\in V_{\rm quad}\cap H\). Then
\(rv^{-1}\in R_{\rm global}\cap H\), with the same class
in \(R_{\rm quad}/V_{\rm quad}\). This proves the assertion
about the arithmetic image.

For completeness, if a closed normal \(L\subseteq N\)
satisfies \(G=LH\), every continuous \(G\)-invariant
homomorphism \(\varphi:L\to E\), with \(E\) finite abelian
and \(\varphi(L\cap H)=0\), extends uniquely by
\[
 \widetilde\varphi(lh)=\varphi(l).
\]
The intersection condition proves well-definedness,
invariance proves multiplicativity, and the compact
quotient map \(L\times H\to G\) proves continuity.
The same description holds for \(N\), since \(G=NH\).
This identifies all finite abelian relative duals and
therefore gives the first isomorphism in
\eqref{eq:uniform-relative-quotients}.

If \(Z\subset L\) are closed normal and \(L/Z\) is pro-odd,
continuous coprime cohomology gives, for every finite
abelian \(2\)-group \(E\),
\[
 H^1(L,E)\xrightarrow{\ \sim\ }H^1(Z,E)^{L/Z}.
\]
Indeed \(H^1(L/Z,E)=H^2(L/Z,E)=0\).
A \(G\)-invariant character on \(Z\) extends uniquely to
a \(G\)-invariant character on \(L\). If it vanishes on
\(Z\cap H\), its restriction to \(L\cap H\) factors through
the pro-odd group \((L\cap H)/(Z\cap H)\), and is zero.
Applying this with \(L=R_{\rm global}\) and
\(Z=V_{\rm global}\) proves the second isomorphism in
\eqref{eq:uniform-relative-quotients}. It does not require
a factorization \(G=V_{\rm global}H\).
The accepted splitting
\(X_2=\operatorname{Hom}(B,\mathbb F_2)\oplus\mathcal S\),
together with the vanishing of \(\mathcal S\) on
\(R_{\rm quad}\), proves the original restriction assertions.

For \(p\in\Pi\), one has \(T_p=Q_pT_p\). Surjectivity of
the first-layer map makes its cokernel zero, so the
hypothesis holds and \(\mathscr K_p\) is injective.
At \(29\) and \(113\), the exact indices
\([T_p:Q_pT_p]=8\) and the three-dimensional first-layer
cokernels prove the natural isomorphism. The established
full arithmetic images \(P_p\), together with
\eqref{eq:uniform-arithmetic-image}, force
\(\ker\mathscr K_p=0\).
At \(163\), the corresponding index is \(4\), the
first-layer cokernel has dimension \(2\), and the
arithmetic image is exactly \(Q_pP_p\). Thus
\(\mathscr K_{163}=0\) on a two-dimensional domain,
and \(163\notin\Pi_{\rm global}\).
\end{proof}

For any one added prime satisfying this criterion, the newly
imposed kernel has infinite index on each preceding normalized
\(R\)- or \(V\)-group. Indeed \(P_p\) has
\(\mathbb Z_2\)-rank \((p-1)/2\): the rational logarithm
identifies its span with the anti-invariant coefficient space,
and \(s-1\) is invertible there because its characteristic
polynomial is \(X^{(p-1)/2}+1\). This also proves that the
simultaneous normalization is a strict, infinite-index
normalization of \(R_{\rm quad}\) and \(V_{\rm quad}\).
If \(\mathscr K_p\) is not injective, the uniform theorem
determines the arithmetic image but does not determine the
full GT image or the relative quadratic quotient.
